\begin{textsample}{POS Dim 1 – 2002 – Score 21.00 – 2002/t000269.txt}  \label{ex:f1_pos_019}
Welcome.
If I could have \textbf{the} first \textbf{slide}, please?
Contrary to calculations made by some engineers, bees can fly, dolphins can swim, and geckos can even climb up \textbf{the} smoothest \textbf{surfaces}.
Now, what I want to do, in \textbf{the} short time I have, is to try to allow each of you to experience \textbf{the} thrill of revealing nature ' s design.
I get to do this all \textbf{the} time, and it ' s just incredible.
I want to try to share just a little bit of that with you in this presentation. \textbf{The} challenge of looking at nature ' s designs — and I ' ll tell you \textbf{the} way that we perceive it, and \textbf{the} way we ' ve used it. \textbf{The} challenge, of course, is to answer this question : what permits this extraordinary performance of animals that allows them basically to go anywhere?
And if we could figure that out, how can we implement those designs?
Well, many biologists will tell engineers, and others, organisms have millions of years to get it right ; they ' re spectacular ; they can do everything wonderfully well.
So, \textbf{the} answer is bio-mimicry : just \textbf{copy} nature directly.
We know from working on animals that \textbf{the} truth is that ' s exactly what you don ' t want to do — because \textbf{evolution} works on \textbf{the} just-good-enough principle, not on a perfecting principle.
And \textbf{the} constraints in building any organism, when you look at it, are really severe.
Natural technologies have incredible constraints.
Think about it.
If you were an engineer and I told you that you had to build an automobile, but it had to start off to be this big, then it had to grow to be full size and had to work every step along \textbf{the} way.
Or think about \textbf{the} fact that if you build an automobile, I ' ll tell you that you also — inside it — have to put a factory that allows you to make another automobile.
And you can absolutely never, absolutely never, because of history and \textbf{the} inherited plan, start with a clean slate.
So, organisms have this important history.
Really \textbf{evolution} works more like a tinkerer than an engineer.
And this is really important when you begin to look at animals.
Instead, we believe you need to be inspired by \textbf{biology}.
You need to discover \textbf{the} general principles of nature, and then use these analogies when they ' re advantageous.
This is a real challenge to do this, because animals, when you start to really look inside them — how they work — appear hopelessly complex.
There ' s no detailed history of \textbf{the} design plans, you can ' t go look it up anywhere.
They have way too many motions for their joints, too many muscles.
Even \textbf{the} simplest animal we think of, something like an insect, and they have more neurons and connections than you can imagine.
How can you make sense of this?
Well, we believed — and we hypothesized — that one way animals could work simply, is if \textbf{the} control of their movements tended to be built into their bodies themselves.
What we discovered was that two-, four-, six- and eight-legged animals all produce \textbf{the} same forces on \textbf{the} ground when they move.
They all work like this kangaroo, they bounce.
And they can be modeled by a spring-mass system that we call \textbf{the} \textbf{spring} mass system because we ' re biomechanists.
It ' s actually a pogo stick.
They all produce \textbf{the} pattern of a pogo stick.
How is that true?
Well, a human, one of your legs works like two legs of a trotting dog, or works like three legs, together as one, of a trotting insect, or four legs as one of a trotting \textbf{crab}.
And then they alternate in their propulsion, but \textbf{the} patterns are all \textbf{the} same.
Almost every organism we ' ve looked at this way — you ' ll see next week, I ' ll give you a hint, there ' ll be an article coming out that says that really big things like T. rex probably couldn ' t do this, but you ' ll see that next week.
Now, what ' s interesting is \textbf{the} animals, then — we said — bounce along \textbf{the} vertical plane this way, and in our collaborations with Pixar, in"A Bug ' s Life,"we discussed \textbf{the} bipedal nature of \textbf{the} characters of \textbf{the} ants.
And we told them, of course, they move in another plane as well.
And they asked us this question.
They say,"Why model just in \textbf{the} sagittal plane or \textbf{the} vertical plane, when you ' re telling us these animals are moving in \textbf{the} \textbf{horizontal} plane?"This is a good question.
Nobody in \textbf{biology} ever modeled it this way.
We took their advice and we modeled \textbf{the} animals moving in \textbf{the} \textbf{horizontal} plane as well.
We took their three legs, we collapsed them down as one.
We got some of \textbf{the} best mathematicians in \textbf{the} world from Princeton to work on this problem.
And we were able to create a model where animals are not only bouncing up and down, but they ' re also bouncing side to side at \textbf{the} same time.
And many organisms fit this kind of pattern.
Now, why is this important to have this model?
Because it ' s very interesting.
When you take this model and you perturb it, you give it a push, as it bumps into something, it self-stabilizes, with no brain or no reflexes, just by \textbf{the} structure alone.
It ' s a beautiful model.
Let ' s look at \textbf{the} mathematics.
That ' s enough! \textbf{The} animals, when you look at them running, appear to be self-stabilizing like this, using basically springy legs.
That is, \textbf{the} legs can do computations on their own ; \textbf{the} control algorithms, in a sense, are embedded in \textbf{the} form of \textbf{the} animal itself.
Why haven ' t we been more inspired by nature and these kinds of discoveries?
Well, I would argue that human technologies are really different from natural technologies, at least they have been so far.
Think about \textbf{the} typical kind of robot that you see.
Human technologies have tended to be large, flat, with right angles, stiff, made of metal.
They have rolling \textbf{devices} and axles.
There are very few motors, very few sensors.
Whereas nature tends to be small, and curved, and it bends and twists, and has legs instead, and appendages, and has many muscles and many, many sensors.
So it ' s a very different design.
However, what ' s changing, what ' s really exciting — and I ' ll show you some of that next — is that as human technology takes on more of \textbf{the} characteristics of nature, then nature really can become a much more useful teacher.
And here ' s one example that ' s really exciting.
This is a collaboration we have with Stanford.
And they developed this new \textbf{technique}, called Shape Deposition Manufacturing.
It ' s a \textbf{technique} where they can mix materials together and mold any shape that they like, and put in \textbf{the} material properties.
They can embed sensors and actuators right in \textbf{the} form itself.
For example, here ' s a leg : \textbf{the} clear part is stiff, \textbf{the} white part is compliant, and you don ' t need any axles there or anything.
It just bends by itself beautifully.
So, you can put those properties in.
It inspired them to show off this design by producing a little robot they named Sprawl.
Our work has also inspired another robot, a biologically inspired bouncing robot, from \textbf{the} University of Michigan and McGill named RHex, for robot hexapod, and this one ' s autonomous.
Let ' s go to \textbf{the} video, and let me show you some of these animals moving and then some of \textbf{the} simple robots that have been inspired by our discoveries.
Here ' s what some of you did this morning, although you did it outside, not on a treadmill.
Here ' s what we do.
This is a death ' s head \textbf{cockroach}.
This is an American \textbf{cockroach} you think you don ' t have in your kitchen.
This is an eight-legged scorpion, six-legged ant, forty-four-legged centipede.
Now, I said all these animals are sort of working like pogo sticks — they ' re bouncing along as they move.
And you can see that in this ghost \textbf{crab}, from \textbf{the} beaches of Panama and North Carolina.
It goes up to four meters per second when it runs.
It actually leaps into \textbf{the} air, and has aerial phases when it does it, like a horse, and you ' ll see it ' s bouncing here.
What we discovered is whether you look at \textbf{the} leg of a human like Richard, or a \textbf{cockroach}, or a \textbf{crab}, or a kangaroo, \textbf{the} \textbf{relative} leg stiffness of that \textbf{spring} is \textbf{the} same for everything we ' ve seen so far.
Now, what good are springy legs then?
What can they do?
Well, we wanted to see if they allowed \textbf{the} animals to have greater stability and maneuverability.
So, we built a terrain that had obstacles three times \textbf{the} hip height of \textbf{the} animals that we ' re looking at.
And we were certain they couldn ' t do this.
And here ' s what they did. \textbf{The} animal ran over it and it didn ' t even slow down!
It didn ' t decrease its preferred speed at all.
We couldn ' t believe that it could do this.
It said to us that if you could build a robot with very simple, springy legs, you could make it as maneuverable as any that ' s ever been built.
Here ' s \textbf{the} first example of that.
This is \textbf{the} Stanford Shape Deposition Manufactured robot, named Sprawl.
It has six legs — there are \textbf{the} tuned, springy legs.
It moves in a gait that an insect uses, and here it is going on \textbf{the} treadmill.
Now, what ' s important about this robot, compared to other robots, is that it can ' t see anything, it can ' t feel anything, it doesn ' t have a brain, yet it can maneuver over these obstacles without any difficulty whatsoever.
It ' s this \textbf{technique} of building \textbf{the} properties into \textbf{the} form.
This is a graduate student.
This is what he ' s doing to his thesis project — very robust, if a graduate student does that to his thesis project.
This is from McGill and University of Michigan.
This is \textbf{the} RHex, making its first outing in a demo.
Same principle : it only has six moving parts, six motors, but it has springy, tuned legs.
It moves in \textbf{the} gait of \textbf{the} insect.
It has \textbf{the} middle leg moving in synchrony with \textbf{the} front, and \textbf{the} hind leg on \textbf{the} other side.
Sort of an alternating tripod, and they can negotiate obstacles just like \textbf{the} animal.
Robert Full : It ' ll go on different \textbf{surfaces} — here ' s sand — although we haven ' t perfected \textbf{the} feet yet, but I ' ll talk about that later.
Here ' s RHex entering \textbf{the} woods.
Again, this robot can ' t see anything, it can ' t feel anything, it has no brain.
It ' s just working with a tuned mechanical system, with very simple parts, but inspired from \textbf{the} fundamental dynamics of \textbf{the} animal.
RF : Here ' s it going down a pathway.
I presented this to \textbf{the} jet propulsion lab at NASA, and they said that they had no ability to go down craters to look for ice, and life, ultimately, on Mars.
And he said — especially with legged-robots, because they ' re way too complicated.
Nothing can do that.
And I talk next.
I showed them this video with \textbf{the} simple design of RHex here.
And just to convince them we should go to Mars in 2011, I tinted \textbf{the} video orange just to give them \textbf{the} sense of being on Mars.
Another reason why animals have extraordinary performance, and can go anywhere, is because they have an effective interaction with \textbf{the} environment. \textbf{The} animal I ' m going to show you, that we studied to look at this, is \textbf{the} gecko.
We have one here and notice its position.
It ' s holding on.
Now I ' m going to challenge you.
I ' m going show you a video.
One of \textbf{the} animals is going to be running on \textbf{the} level, and \textbf{the} other one ' s going to be running up a wall.
Which one ' s which?
They ' re going at a meter a second.
How many think \textbf{the} one on \textbf{the} \textbf{left} is running up \textbf{the} wall? \textbf{Okay}. \textbf{The} point is it ' s really hard to tell, isn ' t it?
It ' s incredible, we looked at students do this and they couldn ' t tell.
They can run up a wall at a meter a second, 15 steps per second, and they look like they ' re running on \textbf{the} level.
How do they do this?
It ' s just phenomenal. \textbf{The} one on \textbf{the} right was going up \textbf{the} hill.
How do they do this?
They have bizarre toes.
They have toes that uncurl like party favors when you blow them out, and then peel off \textbf{the} \textbf{surface}, like tape.
Like if we had a piece of tape now, we ' d peel it this way.
They do this with their toes.
It ' s bizarre!
This peeling inspired iRobot — that we work with — to build Mecho-Geckos.
Here ' s a legged version and a tractor version, or a bulldozer version.
Let ' s see some of \textbf{the} geckos move with some video, and then I ' ll show you a little bit of a \textbf{clip} of \textbf{the} robots.
Here ' s \textbf{the} gecko running up a vertical \textbf{surface}.
There it goes, in real time.
There it goes again.
Obviously, we have to slow this down a little bit.
You can ' t use regular cameras.
You have to take 1, 000 pictures per second to see this.
And here ' s some video at 1, 000 frames per second.
Now, I want you to look at \textbf{the} animal ' s back.
Do you see how much it ' s bending like that?
We can ' t figure that out — that ' s an unsolved mystery.
We don ' t know how it works.
If you have a son or a daughter that wants to come to Berkeley, come to my lab and we ' ll figure this out. \textbf{Okay}, send them to Berkeley because that ' s \textbf{the} next thing I want to do.
Here ' s \textbf{the} gecko mill.
It ' s a see-through treadmill with a see-through treadmill belt, so we can watch \textbf{the} animal ' s feet, and videotape them through \textbf{the} treadmill belt, to see how they move.
Here ' s \textbf{the} animal that we have here, running on a vertical \textbf{surface}.
Pick a foot and try to watch a toe, and see if you can see what \textbf{the} animal ' s doing.
See it uncurl and then peel these toes.
It can do this in 14 milliseconds.
It ' s unbelievable.
Here are \textbf{the} robots that they inspire, \textbf{the} Mecho-Geckos from iRobot.
First we ' ll see \textbf{the} animals toes peeling — look at that.
And here ' s \textbf{the} peeling action of \textbf{the} Mecho-Gecko.
It uses a pressure- sensitive adhesive to do it.
Peeling in \textbf{the} animal.
Peeling in \textbf{the} Mecho-Gecko — that allows them climb autonomously.
Can go on \textbf{the} flat \textbf{surface}, transition to a wall, and then go onto a ceiling.
There ' s \textbf{the} bulldozer version.
Now, it doesn ' t use pressure-sensitive glue. \textbf{The} animal does not use that.
But that ' s what we ' re limited to, at \textbf{the} moment.
What does \textbf{the} animal do? \textbf{The} animal has weird toes.
And if you look at \textbf{the} toes, they have these little leaves there, and if you blow them up and zoom in, you ' ll see that ' s there ' s little striations in these leaves.
And if you zoom in 270 times, you ' ll see it looks like a rug.
And if you blow that up, and zoom in 900 times, you see there are hairs there, tiny hairs.
And if you look carefully, those tiny hairs have striations.
And if you zoom in on those 30, 000 times, you ' ll see each hair has split ends.
And if you blow those up, they have these little structures on \textbf{the} end. \textbf{The} smallest branch of \textbf{the} hairs looks like spatulae, and an animal like that has one billion of these nano-size split ends, to get very close to \textbf{the} \textbf{surface}.
In fact, there ' s \textbf{the} \textbf{diameter} of your hair — a gecko has two million of these, and each hair has 100 to 1, 000 split ends.
Think of \textbf{the} contact of that that ' s possible.
We were fortunate to work with another group at Stanford that built us a special manned sensor, that we were able to measure \textbf{the} force of an individual hair.
Here ' s an individual hair with a little split end there.
When we measured \textbf{the} forces, they were enormous.
They were so large that a patch of hairs about this size — \textbf{the} gecko ' s foot could support \textbf{the} weight of a small child, about 40 \textbf{pounds}, easily.
Now, how do they do it?
We ' ve recently discovered this.
Do they do it by friction?
No, force is too low.
Do they do it by electrostatics?
No, you can change \textbf{the} charge — they still hold on.
Do they do it by interlocking?
That ' s kind of a like a Velcro-like thing.
No, you can put them on molecular smooth \textbf{surfaces} — they don ' t do it.
How about suction?
They stick on in a vacuum.
How about \textbf{wet} adhesion?
Or capillary adhesion?
They don ' t have any glue, and they even stick under water just fine.
If you put their foot under water, they grab on.
How do they do it then?
Believe it or not, they grab on by intermolecular forces, by Van der Waals forces.
You know, you probably had this a long time ago in \textbf{chemistry}, where you had these two atoms, they ' re close together, and \textbf{the} electrons are moving around.
That tiny force is sufficient to allow them to do that because it ' s added up so many times with these small structures.
What we ' re doing is, we ' re taking that inspiration of \textbf{the} hairs, and with another colleague at Berkeley, we ' re manufacturing them.
And just recently we ' ve made a breakthrough, where we now believe we ' re going to be able to create \textbf{the} first synthetic, self-cleaning, dry adhesive.
Many companies are interested in this.
We also presented to Nike even.
We ' ll see where this goes.
We were so excited about this that we realized that that small-size scale — and where everything gets sticky, and gravity doesn ' t matter anymore — we needed to look at ants and their feet, because one of my other colleagues at Berkeley has built a six-millimeter silicone robot with legs.
But it gets stuck.
It doesn ' t move very well.
But \textbf{the} ants do, and we ' ll figure out why, so that ultimately we ' ll make this move.
And imagine : you ' re going to be able to have swarms of these six-millimeter robots available to run around.
Where ' s this going?
I think you can see it already.
Clearly, \textbf{the} Internet is already having eyes and ears, you have \textbf{web} cams and so forth.
But it ' s going to also have legs and hands.
You ' re going to be able to do programmable work through these kinds of robots, so that you can run, fly and swim anywhere.
We saw David Kelly is at \textbf{the} beginning of that with his \textbf{fish}.
So, in conclusion, I think \textbf{the} message is clear.
If you need a message, if nature ' s not enough, if you care about \textbf{search} and rescue, or mine clearance, or medicine, or \textbf{the} various things we ' re working on, we must preserve nature ' s designs, otherwise these secrets will be lost forever.
Thank you.

% matched lemmas: biology, chemistry, clip, cockroach, copy, crab, device, diameter, evolution, fish, horizontal, left, okay, pound, relative, search, slide, spring, surface, technique, the, web, wet
\end{textsample}
